% =====================================================================
%  DBMS Lab Record - Project Report (Library Management System)
%  Continuation of the record: page numbering starts at 104.
%  Build: latexmk -xelatex project.tex
% =====================================================================
% =====================================================================
%  DBMS Lab Record Book
%  Build: latexmk -pdf record.tex
% =====================================================================
\documentclass[12pt,a4paper]{article}

\usepackage{fontspec}
\setmainfont{Liberation Serif}         % metric-compatible with Times New Roman
\setmonofont{Nimbus Mono PS}[Scale=0.87]
\usepackage[a4paper,left=0.75in,right=0.75in,top=1in,bottom=1in]{geometry}
\usepackage{tikz}
\usepackage{eso-pic}
\usepackage{tabularx}
\usepackage{array}
\usepackage{fancyvrb}
\usepackage[normalem]{ulem}

\pagestyle{empty}
\setlength{\parindent}{0pt}
\setlength{\parskip}{0pt}

% ---------- page border (single line, 30pt from the page edge) ----------
\AddToShipoutPictureBG{%
  \begin{tikzpicture}[remember picture,overlay]
    \draw[line width=0.5pt]
      ([xshift=30pt,yshift=-30pt]current page.north west)
      rectangle
      ([xshift=-30pt,yshift=30pt]current page.south east);
  \end{tikzpicture}%
}

% ---------- page number: outside the border, 13pt from the page edge ------
\AddToShipoutPictureFG{%
  \begin{tikzpicture}[remember picture,overlay]
    \node[anchor=south, inner sep=0pt] at ([yshift=13pt]current page.south)
      {\fontsize{12}{14}\selectfont\thepage};
  \end{tikzpicture}%
}

% ---------- headings (sizes taken from DBMS_lab_Format.docx) ----------
\newcommand{\mainheading}[1]{%
  \begin{center}\textbf{\fontsize{18}{22}\selectfont #1}\end{center}\vspace{6pt}}
\newcommand{\subheading}[1]{%
  \textbf{\fontsize{16}{19}\selectfont\uline{#1}}\par\vspace{10pt}}
\newcommand{\subsubheading}[1]{%
  \textbf{\fontsize{14}{17}\selectfont\uline{#1}}\par\vspace{10pt}}
\newcommand{\qheading}[1]{%
  \textbf{\fontsize{14}{17}\selectfont\uline{#1}}\par\vspace{10pt}}
\newcommand{\codeheading}[1]{%
  \textbf{\fontsize{12}{15}\selectfont\uline{#1}}\par\vspace{4pt}}

% ---------- code / output blocks ----------
\DefineVerbatimEnvironment{sqlcode}{Verbatim}%
  {fontsize=\small,formatcom=\ttfamily,xleftmargin=0pt}
\DefineVerbatimEnvironment{sqlout}{Verbatim}%
  {fontsize=\small,formatcom=\ttfamily,xleftmargin=0pt}


\usepackage{graphicx}
\graphicspath{{img/}}

\setcounter{page}{121}

% ---------- report headings (layout taken from the sample report) ----------
\newcommand{\projtitle}[1]{%
  \textbf{\fontsize{22}{26}\selectfont #1}\par\vspace{16pt}}
\newcommand{\projhead}[1]{%
  \textbf{\fontsize{15}{18}\selectfont #1}\par\vspace{8pt}}
\newcommand{\projsub}[1]{%
  \textbf{\fontsize{13}{16}\selectfont #1}\par\vspace{6pt}}
\newcommand{\projitem}[2]{%
  \noindent\textbf{\fontsize{12.5}{15}\selectfont #1}\ \ {\fontsize{12.5}{15}\selectfont #2}\par\vspace{6pt}}
\newcommand{\projtext}[1]{%
  {\fontsize{12}{15.6}\selectfont #1\par}\vspace{4pt}}
\newcommand{\shot}[2][430pt]{%
  \noindent\includegraphics[width=#1]{#2}\par\vspace{16pt}}

\begin{document}

\projtitle{Library Management System}

\projhead{Abstract}

{\fontsize{12}{15.6}\selectfont
\setlength{\parskip}{8pt}

The Library Management System is a simple Microsoft Access based database
application designed to record and manage the day to day activities of a
college library. It keeps track of the books held by the library, the members
who borrow them, and every issue and return transaction, so that the librarian
can work from a single screen instead of a set of registers.

The database consists of five related tables --- Authors, Books, BookIssues,
Members and Publishers --- linked through one-to-many relationships so that the
same author, publisher or member is never repeated inconsistently. A book is
catalogued once and then issued many times, which keeps the issue history
complete without duplicating the book details. Records are entered through a
set of user friendly forms rather than by typing directly into the tables: the
main menu opens the member form, the book form and the issue/return form, and
each form validates and saves its own record.

A query module is used to generate the information the librarian needs --- the
complete book catalogue, the list of books currently available, the active
loans, the overdue books, and the borrowing history of each member. These
queries feed a report that groups the overdue books by the month in which they
fell due and lists the member's name and phone number beside each one, so that
reminders can be sent quickly.

Overall, the project demonstrates the practical use of Microsoft Access for
building a small scale information system that is efficient, easy to use, and
effective for library management.

\par}

\vspace{14pt}
\projhead{Screenshots}

\projsub{Tables}

\projitem{1.}{Books Table}
\shot[470pt]{books_datasheet.jpeg}

\clearpage
\projhead{Forms}

\projitem{1.}{Main Menu Form}
\projtext{The main menu is the starting point of the application. It opens the
book catalogue, the member form, the issue/return form and the overdue books
report from a single window.}
\shot{mainmenu.jpeg}

\projitem{2.}{Members Form}
\projtext{The members form is used to add and edit library members, recording
the name, email, phone number, join date and member type, and whether the
membership is currently active.}
\shot{members_form.jpeg}

\clearpage
\projitem{3.}{Book Issues Form}
\projtext{The book issues form records a loan against a book and a member. It
stores the issue date and the due date, and the return date is filled in when
the book comes back.}
\shot{bookissues_form.jpeg}

\vspace{10pt}
\projhead{Queries}

\projitem{1.}{Book Catalogue Query}
\projtext{The book catalogue query lists every book with its author, publisher,
category and shelf location, so that the shelf can be found without opening the
individual tables.}
\shot{qry_bookctlg.jpeg}

\clearpage
\projhead{Reports}

\projitem{1.}{Overdue Books Report}
\projtext{The overdue books report lists the books that were not returned by
their due date, grouped by the month in which they fell due, with the member's
name and phone number so that a reminder can be sent.}
\shot[480pt]{rpt_overdue.jpeg}

\clearpage
\projhead{Relationships}

\projtext{The relationships window shows how the five tables are joined. Each
author and each publisher is related to many books, each book and each member
is related to many issues, and every relationship enforces referential
integrity.}
\shot[470pt]{relationships.jpeg}

\end{document}
